\section{Заключение}

В рамках лабораторных работ~1 и~2 реализован полный цикл работы
с графом: генерация связного ациклического графа по распределению
Рэлея--Райса с softmax-нормировкой степеней, вычисление.
характеристик (эксцентриситеты, радиус, диаметр, центр), построение
матриц минимальных и максимальных маршрутов длины~$k$ методом
Шимбелла, перечисление всех простых путей между двумя вершинами,
обход рёбер графа в ширину, нахождение кратчайшего пути алгоритмом
Дейкстры с выводом вектора расстояний и сравнение алгоритмов по
числу итераций.

Все пункты заданий выполнены. Граф из лабораторной~1
переиспользуется в лабораторной~2 без перегенерации. Результаты
каждого шага сохраняются в текстовые файлы и передаются
в Python-скрипты для визуализации без перезапуска исполняемого файла.

\subsection{Сильные стороны реализации}

Архитектура разделена на независимые слои: граф не знает о
генераторе, генератор не знает об алгоритмах, алгоритмы не
знают о визуализации. Это позволяет заменять любой компонент
без изменения остальных.

Шаблонный \texttt{GraphBase} даёт одну
кодовую базу для графа и сети за счёт
статического полиморфизма. Шаблонный генератор с лямбдами
переиспользуется для обоих типов графов и любого закона
генерации весов без дублирования алгоритма построения.

\texttt{std::optional} в матрицах Шимбелла и в интерфейсе
графа исключает «специальные» значения и делает семантику
«ребра нет» / «пути нет» явной на уровне типов.

Наличие веб интерфейса с полной визуализацией и анимацией алгоритмов на графах.

\subsection{Слабые стороны реализации}

Метод Шимбелла работает за $O(k \cdot n^3)$ и хранит две
матрицы $n \times n$ в памяти.

Поиск всех простых путей имеет факториальную сложность. На
разреженных ациклических графах это допустимо, но никакой защиты на больших числах не предусмотрено.

Передача данных между C++ и Python реализована через текстовые файлы:
C++ записывает результат в \texttt{assets/txt/}, Python читает его
при запуске скрипта. Формат намеренно простой для учебной задачи.
В production-окружении такое взаимодействие между процессами
заменяется брокерами сообщений (RabbitMQ, Kafka) или очередями задач
(Celery, ZeroMQ)~--- они обеспечивают типизированные сообщения,
подтверждение доставки и горизонтальное масштабирование.

\subsection{Масштабирование}


Замена файлового обмена на брокер сообщений. Это развязяжет процессы по времени ввиду асинхронности, позволит запускать несколько воркеров визуализации параллельно и добавлит гарантии доставки.

Переход на WebSocket, который позволит C++ пушить результат напрямую в браузер:
визуализация появится сразу по завершении вычислений без перезагрузки
страницы. 

Замена Статичные PNG на интерактивные графы через  Plotly --- с возможностью масштабирования, выделения вершин и фильтрации рёбер по весу прямо в браузере.


Параллелизация алгоритмической части программы. Метод Шимбелла считает $n$ независимых строк матрицы на каждом шаге, поиск всех путей естественно разбивается по стартовым рёбрам.

Покрытие тестами \textit{GTest} крайние случаи генератора, корректность матриц Шимбелла и полноту поиска путей.